\section{Conclusion}

In our brief survey covering over 20 years of DPG history, we identified a number of key issues hindering the mainstream adoption of DPGs, despite their theoretically strong advantage of not storing any site-related secrets. 
In this paper, we have chosen to view these problems as opportunities for improvement rather than disqualifying the field of DPG writ large.

MFDPG completely solves some of the major drawbacks of existing DPGs, such as password policy compatibility and multi-factor authentication support, and at a minimum, makes significant progress towards remedying other flaws, such as revocability and brute-force susceptibility. By cryptographically incorporating the entropy of multiple strong authentication factors into the password generation process, it also has the effect of turning any password-based authentication flow into a secure multi-factor authenticated login process.

In light of recent major security incidents with popular password managers, researchers are rightly taking a second look at the way we handle password management. 
While there might be a temptation to ignore DPGs as an outdated approach, we hope this work serves as a launching point for further research into secure, usable DPGs, and believe DPGs should remain an important part of this conversation moving forward.